% Verbatim copy of /Users/alingalatan/workspaces/grothe/handoff/manuscript/legacy/ch31-gpt-averaged.tex (file dated 2026-10-25). Only this line was added.
% Consolidated source module (formerly Chapter 31). Owner: lower-bound/11-history-literature.tex
\section{Lower-bound GPT: averaged profile comparison, the corrected gluing step, and a visible-averaged Laplace window (20 October)}\label{ch:gpt-averaged}
\chapsummary{Visibly averaged comparisons and the corrected gluing step. Failed pointwise or single-scale comparisons remain negative results, not inputs to later proofs.}

\paragraph{Provenance and status (IDE note, 20 October).} Lower-bound GPT's \code{proof/averaged-comparison.tex} from \pathref{gpt-reviews/fibre-domination-20261020/}, ported verbatim apart
from this header and the local macros. \emph{Our replay} \st{C}: all six checkers PASS under \code{python3 -I}: \code{check\_profiles}, \code{check\_profile\_tail\_polynomials},
\code{check\_ch29\_gluing}, \code{check\_fibre\_counterexample}, \code{check\_visible\_laplace\_tail} and \code{check\_adaptive\_constants}. The corrections are applied in
\autoref{ch:forest}.
\providecommand{\EE}{\mathbb E}\providecommand{\CC}{\mathbb C}\providecommand{\eps}{\varepsilon}\providecommand{\HH}{\mathcal H}\providecommand{\Var}{\operatorname{Var}}

\subsection{Scope}
The proposed deterministic-fibre comparison is false. The relevant
comparison averages the actual visible variable, retaining its
correlation with the A/B profiles. The results below give an exact
first-order profile theorem and a finite-window formulation for remote
residuals. They do not prove the global polarized profile assertion or
the finite-amplitude shared A/B comparison. The manuscript's numerical
Grothendieck endpoints are unchanged.

\subsection{First-order ordering of proper Gaussian profiles}
\label{avg:firstorder}

Let \(T\) be exponential of mean one, \(G\) a standard proper complex
Gaussian, \(a,b>0\), and
\[
\ell(t)=\sqrt t(a+bt),\quad
h(v,s)=\EE|s+\sqrt vG|,\quad
w_v(t)=\partial_v[2h(v,\ell(t))]\qquad(v>0).
\]
For a nonnegative polynomial \(p\) with \(\EE p(T)=1\), define
\[
A_p(t)=e^t\int_t^\infty(p(u)-1)e^{-u}\,du.
\]

\begin{theorem}[Exact averaged first variation]
\label{avg:profile}
The function \(w_v\) is strictly decreasing on \([0,\infty)\), and
\[
\left.\frac{d}{ds}\EE\bigl[2h(v+s(p(T)-1),\ell(T))\bigr]
\right|_{s=0+}
=\EE[w_v'(T)A_p(T)].
\]
In particular, if \(A_p\ge0\) and it is not identically zero, this
derivative is strictly negative. For the admissible proper profiles,
\[
\begin{array}{c|c}
p(t)&A_p(t)\\ \hline
(t-1)^2&t^2\\
t&t\\
t^2/2&t^2/2+t
\end{array}
\]
so every nontrivial nonnegative mixture of these profiles has strictly
negative first variation away from the constant profile at fixed
total variance.
\end{theorem}
\begin{proof}
Gaussian heat differentiation and the Laplace formula for inverse
distance give
\[
w_v(t)=\frac12\EE|\ell(t)+\sqrt vG|^{-1}
=\frac1{2\sqrt\pi}\int_0^\infty
\frac{r^{-1/2}}{1+vr}
\exp\left(-\frac{r\ell(t)^2}{1+vr}\right)\,dr.
\]
The derivative of \(\ell(t)^2=t(a+bt)^2\) is strictly positive, so
differentiating this positive integral proves \(w_v'(t)<0\).
For \(v>0\) the differentiated integral is finite even at \(t=0\).

The variation of the variance profile is \(p-1\). For \(0\le s<v\)
its variance is at least \(v-s>0\), justifying differentiation.
Write \(B(t)=\int_t^\infty(p(u)-1)e^{-u}\,du\); then
\(B(0)=0\), \(B'(t)=-(p(t)-1)e^{-t}\), and integration by parts gives
\[
\int_0^\infty w_v(t)(p(t)-1)e^{-t}\,dt
=\int_0^\infty w_v'(t)B(t)\,dt.
\]
Polynomial growth and exponential decay eliminate the boundary terms.
The table follows by direct integration, or the polynomial equation
\(A_p-A_p'=p-1\), with \(A_p(0)=0\).
\end{proof}

At an interior maximizer \(v_*\) of
\(2\EE h(v,\ell(T))-v/\eps\), one has \(\EE w_{v_*}=1/\eps\).
Thus adding the corresponding proper Gaussian sector has exact
first-order penalties
\[
d_A=\EE[-T^2w_{v_*}'(T)],\qquad
d_{Y,W}=\EE[-Tw_{v_*}'(T)],\qquad
d_B=\EE[-(T^2/2+T)w_{v_*}'(T)].
\]
These are analytic identities and strict signs, not new interval
enclosures of \autoref{ch:forest}'s numerical constants \(0.466,0.70,0.93\).
They control a first variation at a constant proper Gaussian
baseline. They do not by themselves prove a finite-amplitude
comparison or control the pseudocovariance of polarized shared
drivers.

\subsubsection{The exact polarized local form}

For the shared-driver perturbation
\[
Q_s=\sqrt vG+\sqrt s\left(\alpha(T-1)X+
                     \beta T\bar X/\sqrt2\right),
\]
where \(G,X\) are independent proper Gaussians, let
\[
k_v(T)=-\frac{\sqrt\pi}{2\sqrt v}e^{-x}I_1(x),
\qquad x=\ell(T)^2/(2v).
\]
Here \(I_1\) is the modified Bessel function. At a Gaussian-optimal
\(v=v_*\), the exact variation of the penalized objective is
\[
\left.\frac{d}{ds}\Lambda(Q_s)\right|_{s=0+}
=-d_A|\alpha|^2-d_B|\beta|^2
 +2\operatorname{Re}(\alpha\beta)K,\qquad
K=\EE[T(T-1)k_v(T)/\sqrt2].
\]
To verify the formula, the perturbation's variance is
\(s(|\alpha|^2(T-1)^2+|\beta|^2T^2/2)\), while its pseudovariance is
\(2s\alpha\beta T(T-1)/\sqrt2\). The trace part of the Gaussian
covariance derivative is \(w_v\); its anisotropic part is
\((F''(\ell)-F'(\ell)/\ell)/4=k_v(T)\), where
\(F(\ell)=2h(v,\ell)\). This gives the displayed expression.

The necessary and sufficient local quadratic condition is
\(d_A\ge0,d_B\ge0,K^2\le d_Ad_B\).
Floating quadrature near the peak gives \(K\approx-0.154995\)
and smallest penalty eigenvalue approximately \(0.418936\).
These two decimals are exploratory, not interval certificates.
The scalar profile theorem does not by itself establish the
determinant condition, or a global finite-amplitude inequality.

\subsubsection{The forbidden profile and its correct interpretation}

For \(p_{\rm bad}(t)=(t-2)^2/2\), the tail polynomial is
\[
A_{p_{\rm bad}}(t)=t^2/2-t,
\]
which changes sign. Hence the theorem does not assert domination
for the profile reported to violate the Gaussian comparison.
This is exactly the needed bidegree-sensitive sanity check.

There is a labeling distinction in \autoref{ch:forest}. The residual
\((T-2)X/\sqrt2\) has this bad profile, but
\[
\frac{(T-2)X}{\sqrt2}
=\frac{(T-1)X}{\sqrt2}-\frac X{\sqrt2}.
\]
It contains a nonzero first-chaos component; it is not a pure
\((2,1)\) cubic. Its first and third-chaos variances are both \(1/2\).
This directly explains the incorrect centering.

A random Gaussian coefficient multiplying
\(H(Z)=Z(T-2)/\sqrt2\) has a different normalized variance profile,
\[
p_D(t)=t(t-2)^2/2,\qquad
A_{p_D}(t)=t^3/2-t^2/2+t
 =t\left((t-\tfrac12)^2/2+\tfrac78\right)>0\quad(t>0).
\]
That random-coefficient construction is also inadmissible in pure
\(H_{2,1}\), but it is \emph{stable} at first order by the theorem.
It must not be identified with the bad profile without the factor \(t\).

\subsection{An exact finite-window comparison after the visible average}
\label{avg:laplace}

Let \(L=Z(a+b|Z|^2)\), and put \(f(u)=2\EE|L+\sqrt u|\).
\autoref{ch:rescub} proves the marginal Gaussian-mixture representation
\(L\stackrel d=\sqrt R G\), with \(R>0\) independent of \(G\).
Equivalently, if \(H(t)=t(a+bt)^2\) and \(h=H^{-1}\), then
\[
e^{-h(u)}=\EE e^{-\lambda u},\qquad R=1/\lambda.
\]

\begin{theorem}[Positive kernel with explicit exponential tail]
Define
\[
\omega(t)=\EE\sqrt{(1-tR)_+},\qquad
d=\frac{4a^3}{27b},\qquad C=\exp(a/(3b)).
\]
Then
\[
f(u)=f(0)+\frac1{\sqrt\pi}\int_0^\infty
 (1-e^{-tu})t^{-3/2}\omega(t)\,dt,\qquad
0\le\omega(t)\le\min(1,Ce^{-dt}).
\]
For every \(Q\) independent of \(Z\), with \(v=\EE|Q|^2<\infty\),
\[
\Lambda(Q)-\Lambda(\sqrt vG_0)
=-\frac1{\sqrt\pi}\int_0^\infty
\left(\EE e^{-t|Q|^2}-\frac1{1+vt}\right)
t^{-3/2}\omega(t)\,dt .
\]
Here \(G_0\) is independent of \(Z\).
\end{theorem}
\begin{proof}
Subtract the \(u=0\) value from the correct fixed-shift Rice formula
\[
f(u)=\frac1{\sqrt\pi}\int_0^\infty
\left[1-\EE\frac{e^{-su/(1+sR)}}{1+sR}\right]s^{-3/2}\,ds.
\]
For each \(R\), set \(t=s/(1+sR)\). The transformed density is
\[
\frac{s^{-3/2}}{1+sR}\,ds
=t^{-3/2}\sqrt{1-tR}\,dt,\qquad 0<t<1/R.
\]
Tonelli applies to the nonnegative increments, proving the kernel.

We establish the critical exponential moment
\(\EE e^{d\lambda}=C\) by nonnegative power series, rather than
assuming analytic continuation of a moment-generating function.
Analyticity of \(h\) at zero and monotone convergence in its Laplace
transform give
\(\EE\lambda^n=(-1)^n(e^{-h})^{(n)}(0)<\infty\).
The formal series \(y(t)=-h(-t)\) solves
\[
y=\frac{t}{(a-by)^2}.
\]
Starting \(y_0=0\), iterate
\(y_{k+1}=t/(a-by_k)^2\). All coefficients are nonnegative;
each fixed coefficient stabilizes at that of the unique formal
solution. For \(0\le t\le d\), the real iterates increase and are
bounded by \(a/(3b)\). Nonnegative-series composition identifies
their real values with their series, and monotone convergence
identifies the limit with the formal solution's sum. At \(t=d\)
the limit solves \(d=y(a-by)^2\) in \([0,a/(3b)]\), hence equals
\(a/(3b)\). The coefficients of \(e^{y(t)}\) are
\(\EE\lambda^n/n!\), so Tonelli gives
\(\EE e^{d\lambda}=\exp(y(d))=C\).
Consequently
\[
\omega(t)\le\mathbb P(\lambda>t)\le Ce^{-dt}.
\]

For independent \(Q\), circularity of \(L\) implies
\(\Lambda(Q)=\EE f(|Q|^2)-v/\eps\), regardless of whether \(Q\)
itself is circular. Subtracting the variance-matched Gaussian
formula proves the last identity. The two increment integrals are
finite: near zero use \(1-e^{-tu}\le tu\), and at infinity use the
exponential kernel bound.
\end{proof}

Define the normalized truncated comparison
\[
\mathcal J_\tau(Q)=\frac1{\sqrt\pi}\int_0^\tau
\left(\EE e^{-t|Q|^2}-\frac1{1+vt}\right)
t^{-3/2}\omega(t)\,dt.
\]
If \(\mathcal J_\tau(Q)\ge-\eta_0\), then
\[
\Lambda(Q)\le S+\eta_0+
\frac{Ce^{-d\tau}}{d\sqrt\pi\,\tau^{3/2}}.
\]
The last term can be divided by \(1+v\tau\) for the one-sided
upper bound, because the negative part of the Laplace difference
is at most \(1/(1+vt)\).
At \(a=699/1250,b=179/1250\), exact rational checks prove tail
bounds \(3.83\cdot10^{-6}\) for \(\tau=50\),
\(4.77\cdot10^{-7}\) for \(\tau=60\), and
\(1.60\cdot10^{-10}\) for \(\tau=100\).

This is an alternative interface for the already controlled remote
problem, not a new cap. It does not treat shared A/B by replacing
the visible marginal with independent mixture variables. The joint
law must satisfy \(|Z|^2=h(|L|^2)\), so the shared visible profiles
remain coupled to the mixture realization.

\subsection{The proposed gluing step does not regularize the law}
\label{avg:glue}

Let \(V\) be all retained drivers and \(W,W'\) independent Gaussian
buffers, independent of \(V\). For every measurable vector-valued \(F\),
\[
\left(V,F\left(V,\frac{W+W'}{\sqrt2}\right)\right)
\stackrel d=(V,F(V,W)).
\]
Thus substituting an averaged Gaussian coordinate into the entire
polynomial preserves the full joint law, not just its variance.
It does not shrink the content, cumulants or payoff deficit.
This corrects the operation written in \autoref{ch:forest}'s Glue bullet.

Averaging independent \emph{polynomial values} is different. For
independent pure circular cubics \(q_j\) in disjoint blocks,
\[
\kappa\left(N^{-1/2}\sum q_j\right)=\kappa(q_1)/N,\qquad
\left\|\Pi_{1,0}\left(q_N|q_N|^2\right)\right\|^2
=\|\Pi_{1,0}(q_1|q_1|^2)\|^2/N^2.
\]
The projection formula follows by expanding \(q_N^2\bar q_N\);
cross-block terms either have zero mean in another block or remain
third chaos. These formulas concern excess kurtosis and polynomial
content, not exact payoff deficit or nonlinear profile content.
They supply no monotonicity of \(\Lambda\).

With shared drivers, even these decay rates can fail. Let
\[
q_N=X A_N,\qquad
A_N=N^{-1/2}\sum_{j=1}^N(|Y_j|^2-1),
\]
where all complex Gaussians are independent. This is pure \(H_{2,1}\).
Centered exponential moments and block independence give
\[
\EE A_N^2=1,\quad \EE A_N^3=2/\sqrt N,\quad
\EE A_N^4=3+6/N.
\]
Hence exactly
\[
\EE|q_N|^2=1,\qquad \kappa(q_N)=4+12/N,\qquad
\Pi_{1,0}(q_N|q_N|^2)=4X/\sqrt N.
\]
The squared content is \(16/N\), and excess kurtosis tends to four.
The limit is \(XG_{\mathbb R}\), a Gaussian of random rank-one
conditional covariance given \(X\). A valid gluing argument must
retain this covariance and pay a quantitative replacement error.

\subsection{Finite exact counterexamples to fibrewise domination}

For \(A(t)=\EE_\theta|1+te^{i\theta}|\) and
\(P(t)=1+t^2/4+t^4/64\), one has for all \(t\ge0\)
\[
P(t)-17t^6/64\le A(t)\le P(t)+t^6.
\]
On \(0\le t\le1\), use the nonnegative series
\(A(t)=\sum_{k\ge0}\binom{1/2}{k}^2t^{2k}\).
The tail coefficient sum is at most \(2-81/64<1\).
For \(t\ge1\), use \(1\le A(t)\le1+t\) and
\(16t^2+t^4\le17t^6\).
Thus for a circular unit-variance \(R\),
\[
\EE|s+R|-\EE|s+G|
\ge\frac{(\EE|R|^4-2)s^2-17\EE|R|^6-384}{64s^5}.
\]
The moment triples \((m_2,m_4,m_6)\) are
\((1,44,11352)\) for \(H(X)\), and \((1,12,540)\) for
\(X^2\bar Y/\sqrt2\).
The proposed two-sign inequality is therefore reversed by at least
\[
\frac{28329}{40000000000}\quad(s=100,\ R=H(X)),\qquad
\frac{169}{268435456}\quad(s=32,\ R=X^2\bar Y/\sqrt2).
\]
These are rational certificates; the smaller-shift numerical
violations in \autoref{ch:forest} remain useful but are not needed here.

\subsection{A finite adaptive-response alternative}

The preceding reassessment proved the exact saddle
\(N_n(\eps)=\max_{U,V}\min_{x,y}\Xi(U,V;x,y)\), where
\[
\Xi=\EE(|U+V-x-y|+|U-V+x-y|)
-\|U\|^2-\|V\|^2/\eps+\|x\|^2+\|y\|^2/\eps.
\]
Here \(U,x\in H_{1,0}\), \(V,y\in H_{2,1}\).
There is a dimension-independent finite algorithm for the inner
minimum. This is a separate route, requiring no Gaussian domination.

In the price norm \(\|(x,y)\|_\eps^2=\|x\|^2+\|y\|^2/\eps\),
replace each modulus by \(\sqrt{|w|^2+\delta^2}\).
Let \(f_\delta,g_\delta\) be the resulting bounded phases,
\(h_\delta=(f_\delta+g_\delta)/2\),
\(k_\delta=(f_\delta-g_\delta)/2\), and
\[
p_\delta(x,y)=(P_{1,0}k_\delta,\eps P_{2,1}h_\delta).
\]
For \(0<\eps\le1\), \(\|p_\delta\|_\eps\le1\),
\(\nabla\Xi_\delta(z)=2(z-p_\delta(z))\);
the function is 2-strongly convex and \(2+2/\delta\)-smooth.

Take \(\delta=1/40000\), \(a_0=b_0=0\), and iterate
\[
c_k=\frac{201a_k+b_k}{202},\quad p_k=p_\delta(c_k),\qquad
a_{k+1}=\frac{40000c_k+p_k}{40001},\quad
b_{k+1}=\frac{200b_k+p_k}{201}.
\]
The self-contained accelerated Lyapunov proof in
\texttt{adaptive-note.md} gives, for every outer \(U,V\),
\[
0\le\Xi(a_{2200})-\min_z\Xi(z)
\le\frac1{20000}+2(200/201)^{2200}<8.434\cdot10^{-5}.
\]
For completeness, its Lyapunov energy is
\(\Xi_\delta(a_k)-\min\Xi_\delta+\|b_k-z_*\|_\eps^2\).
It starts at most two and contracts by \(200/201\): the standard
smooth-descent and strong-convexity inequalities yield this directly
since \(1/201\le1/\sqrt{40001}\).
The extra \(1/20000\) is the uniform smoothing error \(2\delta\).
The companion note supplies the full square expansion.

This counts exact Gaussian projection-oracle calls. Approximate
numerical projections need their own error accounting or a certified
final residual bracket. The outer tensor supremum remains open;
the displayed tolerance is in original-norm units and is not
automatically the same as a row's K allowance.

\subsection{Source integration and status}

The first-order profile theorem can be inserted beside \autoref{ch:forest}'s
numerical table. Replace the Glue scaling claim by the exact
Gaussian-isometry statement and an open quantitative joint
replacement obligation. Correct the bad-profile label by retaining
the factor \(T\) for a random diagonal coefficient.

The finite-amplitude polarized statement in \autoref{ch:forest} is supported
there by 30 floating-point starts. It should remain labeled numerical
evidence unless a separate all-parameter proof or certificate is
provided. The present first-order theorem does not promote that
statement to a theorem.

All supporting exact scripts pass in isolated Python. They check
finite algebra, moment identities and rational constants; the
analytic proofs above carry the general statements. None is a
complete certificate for \(K_G^{\CC}>1.402\), and no live manuscript
or earlier acceptance data were changed.
